
% Compile with XeLaTeX.
\documentclass[aspectratio=43,11pt]{beamer}

\usepackage{fontspec}
\usepackage{microtype}
\usepackage{mathtools}
\usepackage{booktabs}
\usepackage{tabularx}
\usepackage{array}
\usepackage{tikz}
\usepackage{colortbl}
\usepackage{ragged2e}
\usepackage{enumitem}
\usepackage{graphicx}
\usepackage{adjustbox}
\usepackage{multirow}
\usepackage{xcolor}

\usetikzlibrary{arrows.meta,positioning,calc,fit,backgrounds}

% Fonts
\setmainfont{Liberation Sans}
\setsansfont{Liberation Sans}
\setmonofont{DejaVu Sans Mono}

% Theme
\usefonttheme{professionalfonts}
\usetheme[numbering=fraction,progressbar=frametitle]{metropolis}
\metroset{block=fill}
\setbeamertemplate{navigation symbols}{}
\setbeamertemplate{caption}{\raggedright\insertcaption\par}
\setbeamertemplate{itemize items}[circle]
\setbeamertemplate{enumerate items}[default]
\setbeamertemplate{footline}{}

% Colors
\definecolor{Navy}{HTML}{0B1F3A}
\definecolor{Teal}{HTML}{0F766E}
\definecolor{Slate}{HTML}{475569}
\definecolor{SoftBg}{HTML}{F8FAFC}
\definecolor{Accent}{HTML}{2563EB}
\definecolor{WarmGray}{HTML}{64748B}

\setbeamercolor{normal text}{fg=Slate,bg=white}
\setbeamercolor{alerted text}{fg=Teal}
\setbeamercolor{frametitle}{fg=Navy,bg=white}
\setbeamercolor{progress bar}{fg=Teal,bg=SoftBg}
\setbeamercolor{block title}{fg=white,bg=Navy}
\setbeamercolor{block body}{fg=Slate,bg=SoftBg}
\setbeamercolor{block title alerted}{fg=white,bg=Teal}
\setbeamercolor{block body alerted}{fg=Slate,bg=SoftBg}
\setbeamercolor{block title example}{fg=white,bg=Accent}
\setbeamercolor{block body example}{fg=Slate,bg=SoftBg}
\setbeamercolor{title}{fg=Navy}
\setbeamercolor{subtitle}{fg=Slate}
\setbeamercolor{section in toc}{fg=Navy}
\setbeamercolor{structure}{fg=Teal}

% Spacing
\setlist[itemize]{leftmargin=1.15em,itemsep=0.28em,topsep=0.15em}
\setlength{\parskip}{0.25em}
\setlength{\parindent}{0pt}
\RaggedRight
\sloppy

\newcommand{\kw}[1]{\textbf{#1}}

\newcommand{\figph}[3]{%
\begin{center}
\setlength{\fboxsep}{8pt}%
\fcolorbox{WarmGray}{SoftBg}{%
\begin{minipage}[c][#1][c]{\dimexpr #2-2\fboxsep-2\fboxrule\relax}
\centering
{\bfseries #3\par}
\vspace{0.35em}
{\footnotesize Insert figure, screenshot, chart, or diagram here.\par}
{\footnotesize Why it helps: anchors the slide in evidence and keeps the text sparse.\par}
\end{minipage}}%
\end{center}
}

\title{Progress Report}
\subtitle{26 Jun 2026 - 31 Jul 2026}
\author{Soham Kulkarni}
\date{}

\begin{document}

% Cover page as requested.
\begin{frame}[plain]
\begin{tikzpicture}[remember picture,overlay]
    \node at (current page.center) {\includegraphics[height=0.98\paperheight]{cover.png}};
\end{tikzpicture}
\end{frame}


\begin{frame}[shrink=10]{Scope evolution and reporting window}
\begin{columns}[T,onlytextwidth]
    \column{0.54\textwidth}
    \begin{itemize}
        \item Started as a pitch; became a measured recovery study for gfx90c on Fedora 43.
        \item The evidence chain now runs through spoofing, benchmarks, and source builds.
        \item The same workflow is being extended to gfx902 and later upstream work.
    \end{itemize}

    \vspace{0.25em}
    \begin{itemize}
        \item Progress reports: 01-07-26 to 15-07-26, then 15-07-26 to 31-07-26.
        \item Source sequence: blogs \#1 to \#8, with \#1 and \#2 preceding the July window.
    \end{itemize}

    \column{0.46\textwidth}
    \figph{3.2cm}{\linewidth}{Timeline placeholder}
\end{columns}
\end{frame}


\begin{frame}{Introduction}
\begin{columns}[T,onlytextwidth]
    \column{0.56\textwidth}
    \begin{itemize}
        \item Project title: \kw{Generalisable Pipeline for Backporting ROCm to Obsolete AMD Hardware}.
        \item Working context: FOSS United Foundation, Season of Commits.
        \item Core case study: gfx90c on a Ryzen 5650U APU, with gfx902 as the next validation target.
        \item Main shift: from hardware speculation to evidence-backed software recovery.
    \end{itemize}
    \column{0.44\textwidth}
    \figph{4.0cm}{\linewidth}{Native setup screenshot}
\end{columns}
\end{frame}


\begin{frame}{Problem statement}
\begin{columns}[T,onlytextwidth]
    \column{0.57\textwidth}
    \begin{itemize}
        \item Fedora can install ROCm natively, but higher-level packages still fail on gfx90c.
        \item PyTorch and rocm-examples miss required gfx90c artifacts.
        \item GPU spoofing helps validation, but it does not match the real silicon path.
        \item The software stack is fragmented across system ROCm, packaged ROCm, and source builds.
    \end{itemize}
    \column{0.43\textwidth}
    \figph{4.0cm}{\linewidth}{Failure surface placeholder}
\end{columns}
\end{frame}


\begin{frame}{Objectives of the project}
\begin{columns}[T,onlytextwidth]
    \column{0.56\textwidth}
    \begin{itemize}
        \item Restore gfx90c support in a repeatable way.
        \item Measure each recovery stage against a native baseline.
        \item Port the same workflow to gfx902.
        \item Keep the path general enough for future dropped AMD targets.
    \end{itemize}
    \column{0.44\textwidth}
    \figph{4.0cm}{\linewidth}{Objective flowchart placeholder}
\end{columns}
\end{frame}


\begin{frame}{Project plan with timeline}
\begin{columns}[T,onlytextwidth]
    \column{0.55\textwidth}
    \begin{itemize}
        \item 6 months: complete gfx90c port recovery.
        \item 9 months: finish an end-to-end CI/CD pipeline.
        \item Next project: contribute directly to Mojo, LLVM, ROCm, PyTorch, and vLLM.
        \item Trigger: a SOTA RDNA 4 GPU from FOSS United Foundation.
    \end{itemize}
    \column{0.45\textwidth}
    \figph{4.0cm}{\linewidth}{Gantt timeline placeholder}
\end{columns}
\end{frame}


\begin{frame}{Literature review - phase 1}
\begin{columns}[T,onlytextwidth]
    \column{0.55\textwidth}
    \begin{itemize}
        \item Blog \#1 sets the motivation and the gfx90c / Vega 7 target.
        \item Blog \#2 shows native ROCm 6.4.4 on Fedora 43 and the first PyTorch failure.
        \item The first progress report compresses these findings into a working baseline.
    \end{itemize}
    \column{0.45\textwidth}
    \figph{4.0cm}{\linewidth}{Post sequence placeholder}
\end{columns}
\end{frame}

\begin{frame}{Literature review - phase 2}
\begin{columns}[T,onlytextwidth]
    \column{0.55\textwidth}
    \begin{itemize}
        \item Blog \#3 tests GPU spoofing via an HSA override.
        \item Blog \#4 adds AMD first-party benchmarks and FHS packaging context.
        \item The first progress report turns this into benchmark-backed evidence.
    \end{itemize}
    \column{0.45\textwidth}
    \figph{4.0cm}{\linewidth}{Benchmark comparison placeholder}
\end{columns}
\end{frame}

\begin{frame}{Literature review - phase 3}
\begin{columns}[T,onlytextwidth]
    \column{0.55\textwidth}
    \begin{itemize}
        \item Blog \#5 builds GPU programming and execution-model context.
        \item Blog \#6 studies rocm-examples failures on gfx90c.
        \item Blogs \#7 and \#8 move into source builds and the Vega execution / memory model.
    \end{itemize}
    \column{0.45\textwidth}
    \figph{4.0cm}{\linewidth}{Assembly and ISA placeholder}
\end{columns}
\end{frame}


\begin{frame}{Gap in the research, technology, and methodology}
\begin{columns}[T,onlytextwidth]
    \column{0.57\textwidth}
    \begin{itemize}
        \item There was no single reproducible recovery workflow.
        \item Packaged binaries omitted target-specific gfx90c data.
        \item Existing fixes were scattered across posts, issues, and commands.
        \item The missing layer was not just support, but a measurable method.
    \end{itemize}
    \column{0.43\textwidth}
    \figph{4.0cm}{\linewidth}{Sankey or gap map placeholder}
\end{columns}
\end{frame}


\begin{frame}{Description of the proposed solution}
\begin{columns}[T,onlytextwidth]
    \column{0.56\textwidth}
    \begin{itemize}
        \item Use a staged recovery pipeline instead of a single workaround.
        \item Move from native validation to spoofing, mixed stacks, containers, and source builds.
        \item Treat each step as a separate diagnostic layer.
        \item Reuse the same workflow for gfx902 and later targets.
    \end{itemize}
    \column{0.44\textwidth}
    \figph{4.0cm}{\linewidth}{Pipeline diagram placeholder}
\end{columns}
\end{frame}


\begin{frame}{Requirement analysis}
\begin{columns}[T,onlytextwidth]
    \column{0.54\textwidth}
    \begin{itemize}
        \item Fedora 43 host with ROCm 6.4.4.
        \item HIP, rocBLAS, PyTorch, and rocm-examples coverage.
        \item A benchmark set that includes bandwidth, GEMM, and runtime checks.
        \item A second device family for cross-validation.
    \end{itemize}
    \column{0.46\textwidth}
    \figph{4.0cm}{\linewidth}{Requirements matrix placeholder}
\end{columns}
\end{frame}


\begin{frame}{Technology stack}
\begin{columns}[T,onlytextwidth]
    \column{0.54\textwidth}
    \begin{itemize}
        \item Fedora 43, ROCm 6.4.4, HIP, HIPCC.
        \item LLVM, clang, rocm-examples, rocBLAS, PyTorch.
        \item CLBlast, clpeak, Geekbench, and AMD benchmark clients.
        \item Bash, CMake, Git, and source-level patching.
    \end{itemize}
    \column{0.46\textwidth}
    \figph{4.0cm}{\linewidth}{Layered stack diagram placeholder}
\end{columns}
\end{frame}


\begin{frame}{Design}
\begin{columns}[T,onlytextwidth]
    \column{0.55\textwidth}
    \begin{itemize}
        \item Design follows a bottom-up recovery model.
        \item Each layer isolates packaging, build, runtime, or hardware limits.
        \item The same design should explain success on gfx90c and failure on missing paths.
    \end{itemize}
    \column{0.45\textwidth}
    \figph{4.0cm}{\linewidth}{Architecture block diagram placeholder}
\end{columns}
\end{frame}


\begin{frame}{Development and implementation - baseline}
\begin{columns}[T,onlytextwidth]
    \column{0.55\textwidth}
    \begin{itemize}
        \item Verified native ROCm on Fedora 43.
        \item Confirmed HIP execution on gfx90c.
        \item Traced the first PyTorch failure to missing rocBLAS kernels.
        \item Established spoofing as a useful but imperfect baseline.
    \end{itemize}
    \column{0.45\textwidth}
    \figph{4.0cm}{\linewidth}{Command output placeholder}
\end{columns}
\end{frame}

\begin{frame}{Development and implementation - source build path}
\begin{columns}[T,onlytextwidth]
    \column{0.55\textwidth}
    \begin{itemize}
        \item Built rocm-examples from source against ROCm 6.4.4.
        \item Reached the missing \texttt{main\_gfx90c.s} barrier.
        \item Ported \texttt{assembly\_to\_executable} and related build pieces.
        \item Studied how .hip, LLVM IR, and AMDGPU code fit together.
    \end{itemize}
    \column{0.45\textwidth}
    \figph{4.0cm}{\linewidth}{Source build / assembly screenshot placeholder}
\end{columns}
\end{frame}


\begin{frame}{Testing and debugging - failure analysis}
\begin{columns}[T,onlytextwidth]
    \column{0.55\textwidth}
    \begin{itemize}
        \item Classified failures into packaging, metadata, missing Tensile data, and hardware-limited cases.
        \item Compared native, spoofed, and source-built behaviour.
        \item Kept benchmark runs attached to each recovery step.
    \end{itemize}
    \column{0.45\textwidth}
    \figph{4.0cm}{\linewidth}{Error taxonomy / Sankey placeholder}
\end{columns}
\end{frame}

\begin{frame}{Testing and debugging - benchmark comparison}
\begin{columns}[T,onlytextwidth]
    \column{0.55\textwidth}
    \begin{itemize}
        \item Benchmarks covered bandwidth, GEMM, clpeak, Geekbench, and rocBLAS.
        \item Spoofing stayed close on some metrics and diverged on others.
        \item The source-build path clarified which failures were structural.
    \end{itemize}
    \column{0.45\textwidth}
    \figph{4.0cm}{\linewidth}{Native vs spoofed graph placeholder}
\end{columns}
\end{frame}


\begin{frame}{Project outcome}
\begin{columns}[T,onlytextwidth]
    \column{0.55\textwidth}
    \begin{itemize}
        \item gfx90c is now a studied, testable target rather than an unknown.
        \item The runtime path works for basic HIP code.
        \item The remaining blockers sit in packaging, kernels, and generated artifacts.
        \item gfx902 is ready for the same validation sequence.
    \end{itemize}
    \column{0.45\textwidth}
    \figph{4.0cm}{\linewidth}{Outcome dashboard placeholder}
\end{columns}
\end{frame}


\begin{frame}{Conclusion}
\begin{columns}[T,onlytextwidth]
    \column{0.55\textwidth}
    \begin{itemize}
        \item The project moved from a broad idea to a reproducible recovery method.
        \item The strongest evidence points to software and packaging gaps.
        \item The next stage is to finish the gfx90c port and then harden the pipeline.
        \item The later roadmap expands into direct upstream work after new hardware arrives.
    \end{itemize}
    \column{0.45\textwidth}
    \figph{4.0cm}{\linewidth}{Roadmap end-state placeholder}
\end{columns}
\end{frame}


\begin{frame}{References}
\small
\begin{itemize}
    \item AMD vis-\`a-vis GPUs \#1 - \#8.
    \item Progress Report No \#1 and Progress Report No \#2.
    \item AMD ROCm documentation, LLVM source history, and the Vega ISA reference guide.
    \item Fedora packaging notes and issue threads related to gfx90c support.
\end{itemize}
\end{frame}

\end{document}
